
%%%%%%%%%%%%%%%%%%%%%%%%%%%%%%%%%%%%%%%%
%           main body
%%%%%%%%%%%%%%%%%%%%%%%%%%%%%%%%%%%%%%%%

%-----------------------------------
%	TITLE PAGE
%-----------------------------------

\title[数学分析]{\texorpdfstring{\LARGE \bfseries 第九章\quad 数项级数}{第九章 数项级数}} % The short title appears at the bottom of every slide, the full title is only on the title page
\author[无穷乘积] % Your institution as it will appear on the bottom of every slide, maybe shorthand to save space
{\Large \bfseries 无穷乘积}
% \author{Singular Value Decomposition} % Your name
% \date{\today} % Date, can be changed to a custom date
\date{}

%------------------------------------------------

\begin{document}

%------------------------------------------------

\begin{frame}

\titlepage % Print the title page as the first slide

\end{frame}

%------------------------------------------------

% \begin{frame}
% \frametitle{内容提要} % Table of contents slide, comment this block out to remove it
% \tableofcontents % Throughout your presentation, if you choose to use \section{} and \subsection{} commands, these will automatically be printed on this slide as an overview of your presentation
% \end{frame}

%-----------------------------------
%	PRESENTATION SLIDES
%-----------------------------------

%------------------------------------------------

\begin{frame}{引言 | {\smaller[2] $\pi$ 与无穷乘积}}

历史上最早出现的两个关于 $\pi$ 的美丽公式, 都是无穷乘积的形式:

\pause
\vspace{0.5em}

{\larger[1]
\begin{empheq}[box={\fancymath[colback=cyan!30, drop lifted shadow, sharp corners]}]{equation*}
\color{red!80} \frac{\pi}{2} = \frac{2}{1} \cdot \frac{2}{3} \cdot \frac{4}{3} \cdot \frac{4}{5} \cdot \frac{6}{5} \cdot \frac{6}{7} \cdots = \prod_{n=1}^\infty \frac{4n^2}{4n^2 - 1} \quad \text{(Wallis, 1655)}
\end{empheq}
}

\pause
\vspace{0.5em}

{\larger[1]
\begin{empheq}[box={\fancymath[colback=cyan!30, drop lifted shadow, sharp corners]}]{equation*}
\color{red!80} \frac{2}{\pi} = \frac{\sqrt{2}}{2} \cdot \frac{\sqrt{2 + \sqrt{2}}}{2} \cdot \frac{\sqrt{2 + \sqrt{2 + \sqrt{2}}}}{2} \cdots \quad \text{(Viète, 1593)}
\end{empheq}
}

\pause
\vspace{0.5em}

这些``无穷多个数相乘''的公式需要精确的定义和分析工具来理解.

\end{frame}

%------------------------------------------------

\begin{frame}{引言 | {\smaller[2] Euler 乘积公式}}

Riemann $\zeta$ 函数与素数分布有深刻的联系:
\[
\zeta(s) = \sum_{n=1}^\infty \frac{1}{n^s} = \prod_{p ~ \text{prime}} \frac{1}{1 - p^{-s}} = \frac{1}{1-2^{-s}} \cdot \frac{1}{1-3^{-s}} \cdot \frac{1}{1-5^{-s}} \cdots
\]

\pause
\vspace{0.5em}

这一等式称为 \textbf{Euler 乘积公式,} 它将 $\zeta$ 函数表示为无穷乘积. 无穷乘积中每个因子对应一个素数, 揭示了 $\zeta$ 函数与素数分布之间的深刻联系.

\pause
\vspace{0.5em}

以上例子表明无穷乘积在数学中有重要的地位. 接下来我们将给出无穷乘积的严格定义, 并研究它的性质.

\end{frame}

%------------------------------------------------

\section[定义]{无穷乘积的定义}

%------------------------------------------------

\begin{frame}{无穷乘积的定义}

设 $p_1, p_2, \ldots$ 是一个数列, 形式地称它们的``积''为\textbf{无穷乘积}:
\[
\prod_{n=1}^\infty p_n = p_1 \cdot p_2 \cdot p_3 \cdots
\]

\pause
\vspace{0.3em}

类比数项级数, 定义前 $n$ 项的\textbf{部分积}为 $P_n = \prod\limits_{k=1}^n p_k.$

\pause

\begin{Def}[无穷乘积的收敛性]
若部分积序列 $\{P_n\}$ 收敛于\alert{非零}有限实数 $P,$ 则称无穷乘积 $\prod\limits_{n=1}^\infty p_n$ \textbf{收敛,} 并称 $P$ 为其值. 若 $\{P_n\}$ 发散或收敛于零, 则称其\textbf{发散.}
\end{Def}

\pause
\vspace{0.3em}

\begin{attnblock}
与数项级数相比, 要求极限值\alert{非零}. 若 $P_n \to 0,$ 称为\textbf{发散到零}, 是一种特殊的发散情况.
\end{attnblock}

\end{frame}

%------------------------------------------------

\section[典型例子]{典型例子}

%------------------------------------------------

\begin{frame}{收敛的必要条件与基本例子}

\begin{thm}[收敛的必要条件]
若 $\prod\limits_{n=1}^\infty p_n$ 收敛, 则 $\lim\limits_{n \to \infty} p_n = 1.$
\end{thm}

\startproof[bg=yellow, fg=red]
$\lim\limits_{n \to \infty} p_n = \lim\limits_{n \to \infty} \dfrac{P_n}{P_{n-1}} = \dfrac{P}{P} = 1.$
\qed

\pause
\vspace{0.5em}

因此收敛的无穷乘积的通项可以写成 $p_n = 1 + a_n,$ 其中 $a_n \to 0.$

\pause
\vspace{0.5em}

\begin{eg}
$\prod\limits_{n=1}^\infty \left(1 + \dfrac{1}{n}\right)$ 满足必要条件, 但部分积 $P_n = n + 1 \to +\infty,$ 故发散.

$\prod\limits_{n=1}^\infty \left(1 - \dfrac{1}{n+1}\right)$ 的部分积 $P_n = \dfrac{1}{n+1} \to 0,$ 发散到零.
\end{eg}

\end{frame}

%------------------------------------------------

\begin{frame}{Wallis 公式}

\begin{thm}[Wallis 公式, 1655]
\[
\frac{\pi}{2} = \prod_{n=1}^\infty \frac{4n^2}{4n^2 - 1} = \lim_{n \to \infty} \frac{((2n)!!)^2}{(2n-1)!! \cdot (2n+1)!!}.
\]
\end{thm}

\startproof[bg=yellow, fg=red]
部分积 $P_n = \dfrac{((2n)!!)^2}{(2n-1)!! \cdot (2n+1)!!}.$ 令 $\displaystyle I_m = \int_0^{\pi/2} \sin^m x ~ \mathrm{d}x,$ 由递推可得
\[
I_{2n} = \frac{(2n-1)!!}{(2n)!!} \cdot \frac{\pi}{2}, \quad I_{2n+1} = \frac{(2n)!!}{(2n+1)!!},
\]
故 $\dfrac{I_{2n}}{I_{2n+1}} = \dfrac{\pi}{2P_n}.$ 由于 $\sin^{2n+1} x \leqslant \sin^{2n} x \leqslant \sin^{2n-1} x$ 在 $[0, \pi/2]$ 上成立, 有
\[
1 \leqslant \frac{I_{2n}}{I_{2n+1}} \leqslant \frac{I_{2n-1}}{I_{2n+1}} = \frac{2n+1}{2n} \to 1,
\]
由夹逼准则 $\dfrac{I_{2n}}{I_{2n+1}} \to 1,$ 故 $P_n \to \dfrac{\pi}{2}.$
\qed

\end{frame}

%------------------------------------------------

\begin{frame}{Viète 公式}

\begin{thm}[Viète 公式, 1593]
对 $x \neq 0,$ 有
\[
\prod_{n=1}^\infty \cos \frac{x}{2^{n}} = \frac{\sin x}{x}.
\]
令 $x = \pi/2,$ 即得 $\dfrac{2}{\pi} = \dfrac{\sqrt{2}}{2} \cdot \dfrac{\sqrt{2+\sqrt{2}}}{2} \cdot \dfrac{\sqrt{2+\sqrt{2+\sqrt{2}}}}{2} \cdots$
\end{thm}

\startproof[bg=yellow, fg=red]
部分积 $P_n = \prod\limits_{k=1}^n \cos \dfrac{x}{2^k}.$ 反复利用二倍角公式 $\sin 2\theta = 2\sin\theta\cos\theta:$
\[
\sin\frac{x}{2^n} \cdot P_n = \frac{\sin x}{2^n},
\]
故 $P_n = \dfrac{\sin x}{2^n \sin\frac{x}{2^n}} \to \dfrac{\sin x}{x}$ 当 $n \to \infty.$
\qed

\end{frame}

%------------------------------------------------

\section[与级数的关系]{与数项级数的关系}

%------------------------------------------------

\begin{frame}{与数项级数的关系 | {\smaller[2] 充要条件}}

\begin{thm}
无穷乘积 $\prod\limits_{n=1}^\infty p_n$ ($p_n > 0$) 收敛的充要条件是数项级数 $\sum\limits_{n=1}^\infty \ln p_n$ 收敛.
\end{thm}

\startproof[bg=yellow, fg=red]
设部分和 $S_n = \sum\limits_{k=1}^n \ln p_k = \ln P_n.$ 则 $P_n \to P > 0$ $\Longleftrightarrow$ $S_n \to \ln P.$
\qed

\pause
\vspace{0.5em}

\begin{cor}
设 $a_n$ 从某项起恒正或恒负, 则 $\prod\limits_{n=1}^\infty (1 + a_n)$ 收敛 $\Longleftrightarrow$ $\sum\limits_{n=1}^\infty a_n$ 收敛.
\end{cor}

这是因为 $a_n \to 0$ 时 $\dfrac{\ln(1 + a_n)}{a_n} \to 1,$ 由比较判别法的极限形式即得.

\pause
\vspace{0.5em}

由此推论可快速判断: Wallis 乘积中 $a_n = -\dfrac{1}{4n^2},$ $\sum\limits_{n=1}^\infty a_n$ 收敛, 故收敛; Viète 乘积中 $a_n = \cos\dfrac{x}{2^n} - 1 \sim -\dfrac{x^2}{2^{2n+1}},$ $\sum\limits_{n=1}^\infty a_n$ 收敛, 故收敛.

\end{frame}

%------------------------------------------------

\section{总结}

%------------------------------------------------

\begin{frame}{总结}

\begin{itemize}
\item[{\larger[2] \color{red} \ding{43}}] \textbf{定义:} $\prod\limits_{n=1}^\infty p_n := \lim_{n\to\infty} P_n$ (要求极限\alert{非零}).

\pause
\vspace{0.8em}

\item[{\larger[2] \color{red} \ding{43}}] \textbf{必要条件:} $p_n \to 1$ (即 $a_n \to 0,$ 其中 $p_n = 1 + a_n$).

\pause
\vspace{0.8em}

\item[{\larger[2] \color{red} \ding{43}}] \textbf{典型例子:}
\begin{align*}
&\text{Wallis:} \quad \frac{\pi}{2} = \prod\limits_{n=1}^\infty \frac{4n^2}{4n^2-1}; \qquad \text{Viète:} \quad \frac{\sin x}{x} = \prod\limits_{n=1}^\infty \cos \frac{x}{2^n}.
\end{align*}

\pause
\vspace{0.4em}

\item[{\larger[2] \color{red} \ding{43}}] \textbf{与数项级数的关系:} $\prod\limits_{n=1}^\infty (1+a_n)$ 收敛 $\Longleftrightarrow$ $\sum\limits_{n=1}^\infty \ln(1+a_n)$ 收敛; 若 $a_n$ 恒号, 还等价于 $\sum\limits_{n=1}^\infty a_n$ 收敛.
\end{itemize}

\end{frame}

%------------------------------------------------

%------------------------------------------------
% % closing page
% \begin{frame}
% \HUGE{\centerline{\bfseries {\color{gray} The} {\color{zkblue} End}}}

% \pause

% \vspace{1.2em}

% \Large{\centerline{\bfseries {\color{gray}Thank} \quad {\color{zkblue} You}}}

% \end{frame}

%----------------------------------------------------------------------------------------

\end{document}
